% Generated draft: requires mathematical and source review.
\documentclass[11pt]{article}
\usepackage[margin=25mm]{geometry}
\usepackage{fontspec}
\setmainfont{Latin Modern Roman}
\usepackage{amsmath,amssymb,mathtools,unicode-math}
\setmathfont{Latin Modern Math}
\usepackage{enumitem,needspace,xurl,hyperref}
\setlength{\parindent}{0pt}
\setlength{\parskip}{5pt}
\setlength{\emergencystretch}{3em}
\newcommand{\N}{\mathbb{N}}
\newcommand{\Z}{\mathbb{Z}}
\newcommand{\Q}{\mathbb{Q}}
\newcommand{\R}{\mathbb{R}}
\newcommand{\C}{\mathbb{C}}
\newcommand{\F}{\mathbb{F}}
\newcommand{\A}{\mathbb{A}}
\newcommand{\PP}{\mathbb{P}}
\newcommand{\E}{\mathbb{E}}
\newcommand{\eps}{\varepsilon}
\newcommand{\dd}{\,\mathrm{d}}
\begin{document}
\section*{UnsolvedMath ID: 800186}
\textbf{Source problem code:} catalogue entry
\subsection{Definitions and mathematical statement}
For a real parameter \(\alpha>0\), interpret \(d_\alpha\) as the generalized divisor function, defined multiplicatively by
\[
d_\alpha(1)=1,\qquad d_\alpha(p^a)=\frac{\alpha(\alpha+1)\cdots(\alpha+a-1)}{a!}\quad(a\geq1)
\]
for every prime \(p\). Equivalently,
\[
\zeta(s)^\alpha=\sum_{n=1}^{\infty}\frac{d_\alpha(n)}{n^s}\qquad(\operatorname{Re}s>1),
\]
where the power is defined by the Euler product. In particular, \(d_3(n)\) counts ordered triples of positive integers whose product is \(n\).

For each fixed real \(\kappa\in(2,3)\), asymptotically evaluate
\[
S_\kappa(x)=\sum_{\substack{n\in\N\\1\leq n\leq x}}d_\kappa(n)d_3(n+1)
\qquad\text{as }x\to\infty,
\]
with \(x\) real and \(x\geq1\). The task is to determine an explicit leading asymptotic expression, including its leading constant. The supplied statement does not define \(d_\kappa\); the definition above records its standard interpretation. It does not specify a required error bound, lower-order terms, or uniformity as \(\kappa\) varies, so none is imposed here.
\subsection{Short English statement}
For each fixed real \(2<\kappa<3\), determine the asymptotic behavior of the sum of the generalized divisor function \(d_\kappa(n)\) multiplied by the ordinary ternary divisor function at \(n+1\). The generalized divisor notation is interpreted in its standard sense; the requested precision is unspecified.
\subsection{Sources}
\begin{itemize}
\item Local UnsolvedMath registry: permanent ID \(800186\), source problem code \texttt{catalogue entry}, title ``Moments of L-functions --- Moments of the Riemann zeta function --- Problem 38.'' Supplied upstream revision: \texttt{d371ee917ac9647fbb06cd6fc529deec6f3dda82}.
\item Hamieh, Kadiri, Ng, \emph{Moments of L-functions Problem List} (2024), Problem 38. Supplied source URL: \url{https://arxiv.org/abs/2405.17800}. The registry attributes the proposal to Hamieh, Kadiri, Ng (2024).
\item Supplied provenance records access on 2026-10-16 and extraction by \texttt{arxiv\_latex}. These are supplied provenance details, not a claim of independent source inspection. The registry labels the problem ``open''; that status has not been independently verified.
\end{itemize}
\end{document}
